\paragraph{Discussão.}
\textbf{(i) Extração de atributos falha; ajuste fino é essencial.} Com o tronco congelado, a BAcc na CV foi
de apenas $\feCV\%$ e $\feTest\%$ no teste --- \emph{abaixo da MLP} treinada do zero. Os atributos da última
camada de uma rede treinada na ImageNet codificam semântica de objetos naturais, e o domínio aqui é muito
distinto: imagens de microscopia com coloração de May--Grünwald--Giemsa, objeto único centrado e,
sobretudo, imagens de $28\times28$ interpoladas, sem as altas frequências que os filtros iniciais da rede esperam.
Uma cabeça linear sobre esses 512 atributos não consegue reconstruir as pistas finas (granulação, cromatina)
que separam IG, MON e BAS. Já o ajuste fino completo atinge $\ftCV\%$ na CV: a transferência é útil
como \emph{inicialização} --- filtros genéricos de borda/textura e um ponto de partida bem condicionado para
uma rede profunda ---, não como extrator fixo. Dada a diferença de $\approx15$ p.p. entre as estratégias,
muito maior que os desvios entre \emph{folds} ($<1$ p.p.), a escolha de FT pela CV é inequívoca.

\textbf{(ii) Estabilidade e generalização.} O desvio entre \emph{folds} é pequeno ($0{,}36$ p.p. na CV e
$0{,}25$ p.p. no teste), indicando baixa sensibilidade à partição. A BAcc média no teste
($\ftFoldBaccMS\%$) ficou ligeiramente \emph{acima} da estimativa por CV ($93{,}77\%$); como cada
\emph{fold} de validação tem apenas $\approx2.700$ imagens (e $\approx195$ basófilos), a diferença está dentro
do erro de estimação, e não há sinal de sobreajuste à seleção. O comitê dos 5 modelos (média das
probabilidades) acrescenta $+0{,}76$ p.p. de BAcc ($\ftEnsBacc\%$), efeito típico de redução de variância por
\emph{ensembling}, ao custo de $5\times$ mais parâmetros e tempo de inferência.

\textbf{(iii) Estrutura dos erros.} A matriz de confusão do comitê (\cref{fig:tlcm}) mantém
o padrão qualitativo das redes anteriores, porém muito atenuado: IG continua sendo a classe mais difícil
(84,8\% de revocação; confunde-se com MON, BAS e NEU) e MON a segunda (90,5\%, com 6,0\% previstos como IG).
Os basófilos, que eram o ponto fraco da MLP (75,0\%), chegam a 95,9\%, e LIN a 97,9\%. As classes com
assinatura cromática forte (EOS, PLT) ficam próximas de 100\%.

\textbf{(iv) Treinamento ainda não saturado.} A melhor época foi a 10ª--12ª em todos os \emph{folds}, com a
BAcc de validação ainda crescendo lentamente sob o cronograma cosseno; o orçamento de 12 épocas e a resolução
$64\times64$ foram impostos pelo uso exclusivo de CPU ($\approx\ftEpochTime$ s/época). Para referência, o
\emph{benchmark} oficial \cite{yang2021} reporta acurácia de 95,8\% para a ResNet-18 treinada
por 100 épocas (Adam) em $28\times28$ e 96,3\% com entradas redimensionadas para $224\times224$; nosso comitê
(Acc $=\ftEnsAcc\%$) e os modelos individuais ($\ftFoldAccMS\%$) estão nessa faixa, e é razoável esperar
ganhos adicionais com mais épocas e maior resolução de entrada em GPU.
